% Einheit 4 – Ausklammern (bank/terme/e4.jsonl, zusammenbau v1.2)
\einheitenkopf*[e4][Ausklammern]{Ausklammern}
\setcounter{aufgabe}{15}


% test: terme-e4-k2-s10-v2
\lbtest{Kannst du das schon? Dann weiter zu „Terme aufstellen und berechnen“}{
\lbtt{a)}{$5x^2 + 10xy - 15x =$ \lbfeld}
}{a) $5x \cdot (x + 2y - 3)$}


% erkennung: terme-e4-k1-s0-v1, terme-e4-k1-s0-v2, terme-e4-k1-s0-v3, terme-e4-k1-s0-v4
\begin{kbaufgabe}{Erkennen, welcher Faktor in jedem Glied steckt}
\begin{kbblock}
\kbanweisung{Was steckt in jedem Glied von …? Kreise den Faktor ein, der in beiden Gliedern steckt. Klammere nichts aus.}
\item[]\hspace*{-\leftmargin}\lbzelle{3}{\kbteil{a)}{$8x + 12$}{}}\hfill\lbzelle{3}{\kbteil{b)}{$9y^2 + 2y$}{}}\hfill\lbzelleleer{3}\par
\end{kbblock}
\begin{kbblock}
\kbteil{c)}{Was steckt in jedem Glied von $6a^2 + 10a$? Kreise ein, was in beiden Gliedern steckt: Zahl, Variable oder beides. Klammere nichts aus.}{}
\end{kbblock}
\begin{kbblock}
\kbteil{d)}{Was steckt in jedem Glied von $9x - 6 + 15x^2$? Kreise ein, was in allen drei Gliedern steckt. Klammere nichts aus.}{}
\end{kbblock}
\end{kbaufgabe}


% vorstufe: terme-e4-k2-s-1-v1, terme-e4-k2-s-1-v2, terme-e4-k2-s-1-v3, terme-e4-k2-s-1-v4
\begin{kbaufgabe}{Jedes Glied als Produkt mit einem Faktor schreiben}
\begin{kbblock}
\kbteil{a)}{Schreibe jedes Glied von $8x + 20$ als Produkt mit dem Faktor 4.}{}
\kbantwort{$4 \cdot$ \kbfeld $+ 4 \cdot$ \kbfeld}
\end{kbblock}
\begin{kbblock}
\kbteil{b)}{Schreibe jedes Glied von $6a + 14$ als Produkt mit dem Faktor 2.}{}
\kbantwort{$2 \cdot$ \kbfeld $+ 2 \cdot$ \kbfeld}
\end{kbblock}
\begin{kbblock}
\kbteil{c)}{Schreibe jedes Glied von $10y + 15$ als Produkt mit dem Faktor 5.}{}
\kbantwort{$5 \cdot$ \kbfeld $+ 5 \cdot$ \kbfeld}
\end{kbblock}
\begin{kbblock}
\kbteil{d)}{Schreibe jedes Glied von $9b + 6$ als Produkt mit dem Faktor 3.}{}
\kbantwort{$3 \cdot$ \kbfeld $+ 3 \cdot$ \kbfeld}
\end{kbblock}
\end{kbaufgabe}


% vorstufe: terme-e4-k2-s0-v1, terme-e4-k2-s0-v2, terme-e4-k2-s0-v3, terme-e4-k2-s0-v4
\begin{kbaufgabe}{Einen vorgegebenen Faktor ausklammern}
\begin{kbblock}
\item[]\hspace*{-\leftmargin}\lbzelle{1}{\kbteil{a)}{Klammere den Faktor 3 aus: $9x + 21$ \quad $3 \cdot ($ \kbfeld $+$ \kbfeld $)$}{}}\par
\end{kbblock}
\begin{kbblock}
\item[]\hspace*{-\leftmargin}\lbzelle{1}{\kbteil{b)}{Klammere den Faktor 2 aus: $4y + 10$ \quad $2 \cdot ($ \kbfeld $+$ \kbfeld $)$}{}}\par
\end{kbblock}
\begin{kbblock}
\item[]\hspace*{-\leftmargin}\lbzelle{1}{\kbteil{c)}{Klammere den Faktor 5 aus: $10a + 35$ \quad $5 \cdot ($ \kbfeld $+$ \kbfeld $)$}{}}\par
\end{kbblock}
\begin{kbblock}
\item[]\hspace*{-\leftmargin}\lbzelle{1}{\kbteil{d)}{Klammere den Faktor 7 aus: $14b + 21$ \quad $7 \cdot ($ \kbfeld $+$ \kbfeld $)$}{}}\par
\end{kbblock}
\end{kbaufgabe}


% leiter: terme-e4-k2-s1-v1, terme-e4-k2-s1-v2, terme-e4-k2-s2-v1, terme-e4-k2-s3-v1, terme-e4-k2-s4-v1, terme-e4-k2-s5-v1, terme-e4-k2-s6-v1, terme-e4-k2-s9-v1, terme-e4-k2-s10-v1
\begin{kbaufgabe}{Ausklammern}
\begin{kbblock}
\item[]\hspace*{-\leftmargin}\lbzelle{3}{\kbteil{a)}{$5x + 10 =$ \lbfeld}{}}\hfill\lbzelle{3}{\kbteil{b)}{$5x + 25 =$ \lbfeld}{}}\hfill\lbzelleleer{3}\par
\end{kbblock}
\begin{kbblock}
\item[]\hspace*{-\leftmargin}\begin{lbflucht}\makebox[1.6em][r]{c)\hspace{0.2em}} & $x^2 + 5x$ & $=$ & \lbfeld \\[\lbfluchtabstand] \makebox[1.6em][r]{d)\hspace{0.2em}} & $6x^2 + 15x$ & $=$ & \lbfeld \\[\lbfluchtabstand] \makebox[1.6em][r]{e)\hspace{0.2em}} & $6x + 6$ & $=$ & \lbfeld \\[\lbfluchtabstand] \makebox[1.6em][r]{f)\hspace{0.2em}} & $3x^2 + 6x + 15$ & $=$ & \lbfeld\end{lbflucht}\par
\end{kbblock}
\begin{kbblock}
\kbanweisung{Klammere aus und mache die Probe.}
\item[]\hspace*{-\leftmargin}\begin{lbflucht}\makebox[1.6em][r]{g)\hspace{0.2em}} & $6x + 30$ & $=$ & \lbfeld\end{lbflucht}\par
\end{kbblock}
\begin{kbblock}
\kbteil{h)}{Nora soll im Term $5x + 20$ den größten gemeinsamen Faktor ausklammern. Sie rechnet so: \rechnung{5x + 20 &= 5 \cdot (x + 20)} Finde den Fehler und rechne richtig.}{}
\item[]\kbraum{2}
\end{kbblock}
\begin{kbblock}
\kbanweisung{Klammere so weit wie möglich aus und mache die Probe.}
\item[]\hspace*{-\leftmargin}\begin{lbflucht}\makebox[1.6em][r]{i)\hspace{0.2em}} & $3a^2 + 6ab + 9a$ & $=$ & \lbfeld\end{lbflucht}\par
\end{kbblock}
\end{kbaufgabe}


% pflicht: terme-e4-k3-s1-v3, terme-e4-k3-s2-v2, terme-e4-k3-s3-v1, terme-e4-k3-s4-v1
\begin{lbkasten}
\begin{kbaufgabe}{Verstanden?}
\begin{kbblock}
\item[]\lbhinweis{Gemischt – erkenne selbst, was zu tun ist.}
\kbteil{a)}{Max hat in vier Termen ausgeklammert. Welche Ergebnisse können nicht stimmen? Begründe, ohne genau zu rechnen.}{}
\kbfrage{(1) $6a + 18 = 6 \cdot (a + 3)$ \\ (2) $10x - 25 = 5 \cdot (2x + 5)$ \\ (3) $7y + 7 = 7 \cdot y$ \\ (4) $9b - 36 = 9 \cdot (b - 4)$}
\item[]\kbraum{2}
\end{kbblock}
\begin{kbblock}
\kbteil{b)}{Entscheide bei jeder Aussage, ob sie wahr oder falsch ist. Begründe, ohne genau zu rechnen.}{}
\kbfrage{(1) Aus jedem Term mit zwei Gliedern kann man einen Faktor größer als 1 ausklammern. \\ (2) Multipliziert man nach dem Ausklammern wieder aus, erhält man immer den Anfangsterm. \\ (3) Klammert man aus zwei Gliedern aus, stehen in der Klammer nie weniger als zwei Glieder.}
\item[]\kbraum{3}
\end{kbblock}
\begin{kbblock}
\kbteil{c)}{Der Term $6 \cdot (b + 3)$ gibt den Flächeninhalt einer Figur aus zwei Rechtecken in cm² an. Beschreibe die Figur in Worten, ohne sie zu zeichnen.}{}
\item[]\kbraum{2}
\end{kbblock}
\begin{kbblock}
\kbteil{d)}{Ein Sportverein kauft 15 Trikots zu je 28 € und 15 Hosen zu je 17 €. Schreibe die Kosten als Term mit Klammer. Der Verein hat 700 €. Reicht das Geld?}{}
\item[]\kbraum{2}
\end{kbblock}
\end{kbaufgabe}
\end{lbkasten}
